\documentclass[10pt,letterpaper]{article}
\usepackage{cornell-notes}

\title{Clause 26.07.32.03: Class Template Cartesian Product View Iterator}
\author{}
\date{}
\setDocTitle{Clause 26.07.32.03: Class Template Cartesian Product View Iterator}
\setDocSubtitle{C++ 2024}
\setDocOwner{C++ 2024 Cornell Notes}
\setDocDate{\today}
\setCornellCollection{C++ 2024 Cornell Notes}
\setCornellUnitType{Clause}
\setCornellUnitNumber{26.07.32.03}
\setCornellUnitTitle{Class Template Cartesian Product View Iterator}
\hypersetup{pdftitle={Clause 26.07.32.03: Class Template Cartesian Product View Iterator},pdfsubject={Cornell notes},pdfauthor={}}

\begin{document}
\maketitle
\begin{CornellOverviewBox}{Overview}
\textbf{Classification:} Standard library - Normative\\[2pt]
\textbf{Learning objective:} Understand the rules, terminology, and semantic consequences associated with Class template cartesian\_product\_view::iterator.
\end{CornellOverviewBox}\begin{CornellSummaryBox}{Summary}
{\sffamily\bfseries\color{Accent} Summary}\par\vspace{3pt}Section 26.7.32.3 focuses on Class template cartesian\_product\_view::iterator. Apply the specified interface together with its mandates, effects, result conditions, exception behavior, and complexity constraints. When applying it, preserve the distinction between mandatory requirements, recommendations, permissions, and behavior deliberately left undefined, unspecified, or implementation-defined.
\end{CornellSummaryBox}

\end{document}
